\begin{frame}{Postopno odkrivanje v okolju align}
    % Naloga 3.2:
    % Za pravilno postopno odkrivanje morate na enem mestu uporabiti ukaz \only,
    % na dveh pa ukaz \onslide.

    Dokaži \emph{binomsko formulo}: za vsaki realni števili \(a\) in \(b\)
    in za vsako naravno število \(n\) velja

    \begin{align*}
    (a+b)^n \only<1>{&= \ldots \\}
    \onslide<2->{&= (a+b)(a+b)\dots(a+b) \\}
    \onslide<3->{&= a^n + n a^{n-1}b + \dots
    + \binom{n}{k}a^{n-k}b^k + \dots + nab^{n-1} + b^n \\}
    \onslide<4->{&= \sum_{k=0}^n \binom{n}{k}a^{n-k}b^k}
    \end{align*}
\end{frame}
